\chapter{Calendar Spreads and Curve Trades}\label{ch:s1:calendar-spreads-and-curve-trades}

On 20 April 2020 the May WTI crude oil contract opened at \$17.73 a barrel and settled at minus \$37.63, the first negative price in its 37 years; the
CFTC's staff found an oversupplied market, a collapse in demand and a delivery point, Cushing in Oklahoma, whose working storage had been near capacity
since March. A calendar spread, long one delivery and short the next, looks like the safest trade in futures: the two legs share almost all their risk
and the spread moves slowly. On the EIA's thirty-nine years of WTI settlements a rule that bets on the spread's mean reversion earned a Sharpe ratio of
0.30 before costs and 0.05 after, from 1986 to 2019; in 2020 it lost 43\% of its capital at a 10\% volatility target, most of it in March, as storage
filled. The build is \texttt{firm.curvestrat}.

\section{Trading the shape of a curve}

A futures curve has a level, a slope and, further out, a curvature. An outright position trades the level; a calendar spread (Book~1, chapter~19)
trades the slope between two deliveries and is almost immune to the level.

\begin{definition}[Curve trade]\label{def:s1:calendar-spreads-and-curve-trades:curve}
A \emph{curve trade}\index{curve trade} is a position in two or more deliveries of the same futures contract (or points of a yield or volatility curve)
whose weights cancel exposure to the curve's level, so that its P\&L comes from changes in the curve's shape: a calendar spread for the slope between
two deliveries, a butterfly for the curvature among three.
\end{definition}

\begin{proposition}[A calendar spread trades the change in carry]\label{prop:s1:calendar-spreads-and-curve-trades:carry}
If the log futures price for delivery $T$ is $\ln F_t(T) = x_t - c_t (T - t) + g(T)$, with $x_t$ a common level, $c_t$ the carry and $g$ a seasonal
component fixed by the delivery date, then a position long delivery $T_a$ and short delivery $T_b > T_a$ has a log P\&L of
\[
d\big(\ln F_t(T_a) - \ln F_t(T_b)\big) = (T_b - T_a)\, dc_t .
\]
\end{proposition}
\begin{proof}
The level $x_t$ and the seasonal terms cancel or are constant; $d[-c_t(T_a - t) + c_t(T_b - t)] = (T_b - T_a)\,dc_t$ since the terms in $c_t\,dt$
cancel.
\end{proof}

A long spread (near minus far) profits when the curve moves toward backwardation and loses when contango deepens; its size scales with the distance
between deliveries. What moves the slope in commodities is inventory: Gorton, Hayashi and Rouwenhorst found that the convenience yield falls, nonlinearly,
as inventories rise, and that the basis and past returns reflect the state of inventories and forecast risk premiums. Low inventories make the curve
backwardated and steep; high inventories push it into contango up to the cost of storage, the full carry of Book~3, chapter~10.

\section{Mean reversion of spreads}

\begin{definition}[Spread mean reversion]\label{def:s1:calendar-spreads-and-curve-trades:mr}
\emph{Spread mean reversion}\index{spread mean reversion} is the tendency of a spread between related prices to return toward a normal level after a
deviation; a spread mean-reversion strategy sells the spread when it is high relative to its recent mean and buys it when it is low.
\end{definition}

Trading a WTI spread over decades needs contract identity. The EIA publishes generic series: contract~1 is whichever expires next, and on the day after
an expiry it becomes the old contract~2. \Cref{lst:s1:calendar-spreads-and-curve-trades:pair} numbers the contracts by serial, holds the nearest pair clear
of expiry (1--2, moving to 2--3 five trading days before the nearby expires), and follows the same two contracts from one day to the next, so that a roll
of the generic series is not mistaken for a price move. The log spread $\ln(F_{\text{near}}/F_{\text{far}})$ of the held pair is persistent but
mean-reverting: its daily autocorrelation is 0.949, a half-life of 13.3 trading days.

The rule (\cref{lst:s1:calendar-spreads-and-curve-trades:band}) computes the log spread's z-score against its last 20 days, enters against it beyond 1.5
(long the spread when it is unusually low, short when high), and leaves when the z-score comes back inside 0.5 or changes sign. The position is sized
at entry to 10\% annual volatility from an EWMA forecast of the spread's dollar P\&L, and each unit traded costs one cent a barrel (assumed).

\begin{center}
\small
\begin{tabular}{@{}lcccc@{}}
\toprule
storage filter: no long spread when carry is below & none & $-80\%$ & $-40\%$ & $-20\%$\\
\midrule
Sharpe ratio before costs, 1986--2019 & 0.30 & 0.33 & 0.31 & 0.39\\
Sharpe ratio after costs, 1986--2019 & 0.05 & 0.08 & 0.05 & 0.14\\
share of days with a position & 47\% & 47\% & 45\% & 41\%\\
2020: March & $-58.3\%$ & $-20.8\%$ & $-10.5\%$ & $-1.5\%$\\
2020: April & $+22.9\%$ & 0.0\% & 0.0\% & 0.0\%\\
2020: the year & $-42.7\%$ & $-28.1\%$ & $-17.8\%$ & $-8.8\%$\\
\bottomrule
\end{tabular}
\end{center}

The mean reversion is real but small: a Sharpe ratio of 0.30 before costs over 34 years, and one cent a barrel per unit traded takes most of it
(\cref{fig:s1:calendar-spreads-and-curve-trades:cumulative}). The record is also not stable: after costs the rule had a Sharpe ratio of 0.47 from
1986 to 2004 and $-0.32$ from 2005 to 2019, so it had been losing for fifteen years before storage ran out. The spread is a slow, cheap-looking trade
whose edge is the size of its costs.

\begin{omfigure}
\begin{tikzpicture}
\begin{axis}[width=11cm,height=5.2cm,xlabel={\small year},ylabel={\small cumulative return (\%)},tick label style={font=\footnotesize},
  xmin=1986,xmax=2021,ymin=-40,ymax=90,grid=major,grid style={gray!25},legend columns=2,legend style={font=\scriptsize,
  at={(0.5,-0.3)},anchor=north,draw=none,fill=none,/tikz/every even column/.append style={column sep=8pt}},
  table/col sep=comma,/pgf/number format/1000 sep={}]
\addplot[omDef,thick] table[x=year,y=none]{figdata/strategies-1/21-calendar-spreads-and-curve-trades/cumulative.csv};
\addplot[omThm,thick] table[x=year,y=filter]{figdata/strategies-1/21-calendar-spreads-and-curve-trades/cumulative.csv};
\legend{no filter,no long spread when carry is below $-20\%$}
\end{axis}
\end{tikzpicture}

\omcaption{The WTI calendar-spread mean-reversion rule after costs, 1986--2020, at a 10\% volatility target: cumulative sum of daily returns, without a
storage filter and with the filter that turns out best in hindsight. Data: \texttt{s1\_curve.book}.}\label{fig:s1:calendar-spreads-and-curve-trades:cumulative}
\end{omfigure}

\section{Storage limits and the 2020 negative price}

\begin{definition}[Storage-limit risk]\label{def:s1:calendar-spreads-and-curve-trades:storage}
\emph{Storage-limit risk}\index{storage-limit risk} is the risk that a commodity's contango widens beyond the cost of storage when storage runs out, so
that a position betting on the spread narrowing (long the nearer delivery, short the farther) loses without the bound that full carry normally
provides.
\end{definition}

Contango is normally bounded: when the far contract is more expensive than the near one plus the cost of storing oil until then, a trader with storage
buys the near, stores it and sells the far. When storage is full, nobody can do that trade, and nothing stops the spread. That is what the rule met.
The pair's carry was $+4.5\%$ a year in early January 2020; the rule went long the spread on 10 January as it fell below its recent mean, and stayed
long, on and off, until 23 April (\cref{fig:s1:calendar-spreads-and-curve-trades:carry2020}). By 31 March the carry was $-216\%$ a year; on 20 April,
$-302\%$; on 21 April, the day after the May contract settled at $-\$37.63$, the June--July pair's carry was $-575\%$. The rule lost 58.3\% of capital in
March and made back 22.9\% in April, as it re-entered at smaller sizes and the spread partly recovered.

\begin{omfigure}
\begin{tikzpicture}
\begin{axis}[width=11cm,height=5.2cm,xlabel={\small trading day of 2020},ylabel={\small carry of the held pair (\% a year)},
  tick label style={font=\footnotesize},xmin=1,xmax=124,ymin=-600,ymax=50,grid=major,grid style={gray!25},
  xtick={1,41,63,84,104},xticklabels={2 Jan,2 Mar,1 Apr,1 May,1 Jun},legend pos=south west,legend style={font=\scriptsize},
  unbounded coords=jump,table/col sep=comma]
\addplot[omDef,thick] table[x=n,y=carry_pct]{figdata/strategies-1/21-calendar-spreads-and-curve-trades/carry2020.csv};
\addplot[only marks,mark=*,mark size=1.1pt,omThm] table[x=n,y=long]{figdata/strategies-1/21-calendar-spreads-and-curve-trades/carry2020.csv};
\addplot[black!50,dashed] coordinates {(1,-20) (124,-20)};
\legend{carry of the pair held,days long the spread,$-20\%$ filter}
\end{axis}
\end{tikzpicture}

\omcaption{The annualised carry, $12\ln(F_{\text{near}}/F_{\text{far}})$, of the WTI pair the rule held in the first half of 2020, from the EIA's
settlements, and the days on which the unfiltered rule was long the spread. Data: \texttt{s1\_curve.wti}.}\label{fig:s1:calendar-spreads-and-curve-trades:carry2020}
\end{omfigure}

A filter that forbids and closes long spreads in very steep contango cuts the loss: to 28.1\% with a threshold of $-80\%$ a year, 17.8\% with $-40\%$,
8.8\% with $-20\%$, which also has the best Sharpe ratio before 2020 (0.14 after costs). It would be easy to present the $-20\%$ filter as the
strategy. It was chosen after seeing 2020, among four candidates, on a sample with one storage crisis. Before 2020 the carry was below $-20\%$ on
13.1\% of days and below $-40\%$ on 3.5\%. The defensible version of the filter comes from the mechanism and not from the backtest: measure how close
storage is to full (inventories, the cost of storage) and stop betting against contango when it is.

\section{Curve trades in the synthetic universe, rates and volatility}

On \texttt{firm.synthfut}'s ten commodities the proposition holds exactly: the log spread between the nearby and the contract twelve months out moves
one for one with the planted carry (a correlation of 1.00), because the model's curves have no noise of their own. The carry reverts to each market's
mean with a half-life of a year, so a spread book against the spread's 252-day z-score has a Sharpe ratio of 1.00 before costs traded daily; after 4
basis points per leg per unit traded it loses ($-0.29$). Rebalanced monthly, it keeps 0.76 of its 1.05 before costs. A slow signal traded daily pays for
noise in its own position.

The same structure appears in other curves. In interest rates, \omterm{def:s1:calendar-spreads-and-curve-trades:curve}{curve trades} are steepeners and flatteners between two maturities, or butterflies,
weighted to be neutral in duration (Book~2); their P\&L is the change in the curve's slope, as here. In volatility, VIX futures are usually in contango,
because they carry a premium over the expected spot VIX. Simon and Campasano found that from 2006 to 2011 the VIX futures basis did not forecast
changes in spot VIX but did forecast changes in VIX futures prices, and that shorting VIX futures in contango and buying them in backwardation, hedged
with mini-S\&P 500 futures, was highly profitable and robust to costs. The volatility curve's storage limit is a crash: the contango that pays the short
is exactly what reverses when volatility spikes.

\section{Strategy files}

\begin{strategyfile}{Front--second spread mean reversion}\label{strat:s1:calendar-spreads-and-curve-trades:front}
\sfield{Who pays you, and why} Hedgers and index rolls that push the near spread away from its normal level for a few weeks.
\sfield{Instruments and venues} Calendar spreads in liquid commodity futures (WTI, Brent, products, grains), traded as listed spreads.
\sfield{Signal} The log spread's z-score against its recent mean (20 days).
\sfield{Sizing and execution} Enter beyond 1.5, leave inside 0.5; sized at entry to a volatility target; the pair clear of expiry.
\sfield{Costs} The listed spread's bid--ask, a tick or more; they are the size of the edge.
\sfield{How it dies} Storage limits, when contango widens without bound; costs.
\sfield{Horizon, capacity, infrastructure} Weeks; capacity in the front spreads is large; contract-identity data.
\sfield{Backtest honestly} Contract identity through expiries; spread quotes and costs; the storage crises in the sample.
\sfield{Sources} CFTC interim staff report (2020); this chapter: 0.30 before costs and 0.05 after, 1986--2019; $-42.7\%$ in 2020.
\end{strategyfile}

\begin{strategyfile}{Curve steepener in commodity futures}\label{strat:s1:calendar-spreads-and-curve-trades:steepener}
\sfield{Who pays you, and why} The inventory cycle: the slope reverts as stocks are drawn or rebuilt.
\sfield{Instruments and venues} Near and twelve-month commodity futures.
\sfield{Signal} The slope against its one-year mean; inventory data.
\sfield{Sizing and execution} Long the spread when contango is unusually deep and inventories are falling, short when backwardation is unusually steep;
rebalanced monthly.
\sfield{Costs} Two legs, two rolls.
\sfield{How it dies} Inventories that keep building; storage limits.
\sfield{Horizon, capacity, infrastructure} Months.
\sfield{Backtest honestly} Inventory data as released, with revisions; the seasonal pattern of the curve separated from its slope.
\sfield{Sources} Gorton, Hayashi and Rouwenhorst (2013); this chapter: 0.76 on the synthetic universe, monthly.
\end{strategyfile}

\begin{strategyfile}{Roll-period calendar spread}\label{strat:s1:calendar-spreads-and-curve-trades:roll}
\sfield{Who pays you, and why} \omterm{def:s1:index-rebalancing:fund}{Index funds} and ETFs that roll large positions on published schedules, pressing the spread they roll.
\sfield{Instruments and venues} The front spreads of index commodities during the roll windows.
\sfield{Signal} The calendar of index rolls (chapter~9's flows applied to futures).
\sfield{Sizing and execution} Take the other side of the roll before it, unwind during it.
\sfield{Costs} Low; competition from other front-runners.
\sfield{How it dies} \omterm{def:s1:index-rebalancing:fund}{Index funds} spreading their rolls over more days or dates.
\sfield{Horizon, capacity, infrastructure} Days.
\sfield{Backtest honestly} The index rules and roll dates in force at each date.
\sfield{Sources} No performance figure verified.
\end{strategyfile}

\begin{strategyfile}{Volatility futures curve trade}\label{strat:s1:calendar-spreads-and-curve-trades:vix}
\sfield{Who pays you, and why} Buyers of volatility protection who pay a premium over the expected VIX.
\sfield{Instruments and venues} VIX futures, hedged with equity index futures.
\sfield{Signal} The basis: futures over spot VIX (contango) or under it (backwardation).
\sfield{Sizing and execution} Short in contango, long in backwardation, equity exposure hedged; small against crash risk.
\sfield{Costs} Moderate; rolls.
\sfield{How it dies} Volatility spikes that reverse contango in a day.
\sfield{Horizon, capacity, infrastructure} Weeks.
\sfield{Backtest honestly} Futures settlements, not the spot index; hedge ratios estimated out of sample; spikes in the sample.
\sfield{Sources} Simon and Campasano (2014).
\end{strategyfile}

\section{Tutorial: minus thirty-seven}\label{tut:s1:calendar-spreads-and-curve-trades:tut}

\begin{tutorial}
\textbf{Goal.} Build the WTI calendar spread by contract identity from the EIA's generic series, trade its mean reversion with and without a storage
filter, follow it through 2020, and repeat the trade on the synthetic curves. \textbf{End state:} the table and the two figures.
\begin{tutsteps}
\item \textbf{Contract identity}: serial numbers, the held pair and its P\&L through expiries.
\omcode{code/firm/curvestrat/firm_curvestrat.py}{26}{59}{Calendar-spread P\&L by contract identity.}\label{lst:s1:calendar-spreads-and-curve-trades:pair}
\item \textbf{The rule}: enter beyond a z-score, leave inside a smaller one, with a filter on long spreads.
\omcode{code/firm/curvestrat/firm_curvestrat.py}{84}{97}{The band rule with a storage filter.}\label{lst:s1:calendar-spreads-and-curve-trades:band}
\item \textbf{Run} \texttt{persistence()}, \texttt{book(threshold)} for the four filters, \texttt{synthetic(252, monthly)} both ways, and
\texttt{fig\_curve.py}.
\end{tutsteps}
\textbf{What to change next.} Replace the carry filter with EIA inventory data at Cushing; add a butterfly on contracts 1, 2 and 3; charge the listed
spread's quoted bid--ask instead of a fixed cent.
\end{tutorial}

\section{Build: calendar spreads}\label{bld:s1:calendar-spreads-and-curve-trades:curvestrat}

\begin{build}
\bfield{Purpose} Calendar spreads by contract identity, their P\&L through expiries, mean-reversion signals and a regime filter.
\bfield{Interface} \texttt{serials(expiry)}, \texttt{pair\_pnl(C, m, near, days\_before, expiry)}, \texttt{zscore(x, window)}, \texttt{positions(z, vol,
target, cap, block)}, \texttt{band(z, vol, enter, leave, target, block)}, \texttt{spread\_carry(log\_spread, days)}.
\bfield{Rules} A contract keeps its serial number through expiries; the pair moves out before the nearby's last days; signals use data up to $t$.
\bfield{Acceptance tests} \texttt{code/firm/curvestrat/tests/}: a hand-made curve through an expiry, with the pair's P\&L by hand; the z-score, the band's
entry and exit and the filter; carry from a log spread.
\bfield{Stretch} Butterflies; inventory filters; listed-spread quotes.
\end{build}

\begin{omsources}
\item CFTC Division of Market Oversight and Office of the Chief Economist, \emph{Interim Staff Report on Trading in NYMEX WTI Crude Oil Futures Contract
Leading up to, on, and around April 20, 2020}, November 2020.
\item G.~B.~Gorton, F.~Hayashi and K.~G.~Rouwenhorst, ``The fundamentals of commodity futures returns'', \emph{Review of Finance} 17(1), 2013.
\item D.~P.~Simon and J.~Campasano, ``The VIX futures basis: evidence and trading strategies'', \emph{Journal of Derivatives} 21(3), 2014.
\item US Energy Information Administration, NYMEX WTI crude oil futures daily settlements, contracts 1 to 4.
\end{omsources}

\section{Exercises}

\begin{exercise}[$\star$]\label{exo:s1:calendar-spreads-and-curve-trades:1}
A log spread has a daily autocorrelation of 0.949. What is its half-life in trading days?
\end{exercise}

\begin{exercise}[$\star$]\label{exo:s1:calendar-spreads-and-curve-trades:2}
Two contracts 21 trading days apart have a log spread of $-0.01$. What is the annualised carry? On 21 April 2020 the held pair's carry was $-575.5\%$ a
year: what was the ratio of the near price to the far one?
\end{exercise}

\begin{exercise}[$\star$]\label{exo:s1:calendar-spreads-and-curve-trades:3}
Why does the generic series ``contract 1 minus contract 2'' mislead a spread backtest on the day after an expiry?
\end{exercise}

\begin{exercise}[$\star\star$]\label{exo:s1:calendar-spreads-and-curve-trades:4}
The unfiltered rule has Sharpe ratios of 0.30 before costs and 0.05 after, at about 10\% volatility. Roughly how much do costs take each year?
\end{exercise}

\begin{exercise}[$\star\star$]\label{exo:s1:calendar-spreads-and-curve-trades:5}
Why did the rule lose in March 2020 and gain in April?
\end{exercise}

\begin{exercise}[$\star\star$]\label{exo:s1:calendar-spreads-and-curve-trades:6}
Why is the $-20\%$ filter's result not evidence that it works?
\end{exercise}

\begin{exercise}[$\star\star\star$]\label{exo:s1:calendar-spreads-and-curve-trades:7}
\emph{Coding.} Run \texttt{synthetic} daily and monthly. Why does the same signal lose after costs daily and earn 0.76 monthly?
\end{exercise}

\begin{exercise}[$\star\star\star$]\label{exo:s1:calendar-spreads-and-curve-trades:8}
\emph{Find the flaw.} ``The front spread has mean-reverted every month for twenty years; we can run it at ten times leverage because the legs hedge
each other.''
\end{exercise}

\section{Problem: Minus Thirty-Seven}

\begin{problem}[{Weekend problem --- a spread meets a storage limit}]\label{pb:s1:calendar-spreads-and-curve-trades:1}
The EIA's WTI settlements, the synthetic curves and the public record.

\medskip\noindent\textbf{Part I --- Curves.}
\begin{enumerate}
\item Define a \omterm{def:s1:calendar-spreads-and-curve-trades:curve}{curve trade} and \omterm{def:s1:calendar-spreads-and-curve-trades:mr}{spread mean reversion}.
\item State and prove the proposition on calendar spreads and carry.
\item What did Gorton, Hayashi and Rouwenhorst find about inventories?
\item What bounds contango, and when does the bound fail?
\end{enumerate}

\medskip\noindent\textbf{Part II --- The WTI spread.}
\begin{enumerate}[resume]
\item How is the spread followed by contract identity?
\item How persistent is the log spread?
\item Describe the rule and its costs.
\item Give its Sharpe ratios before and after costs, 1986--2019.
\end{enumerate}

\medskip\noindent\textbf{Part III --- 2020.}
\begin{enumerate}[resume]
\item What did the CFTC's staff report find about 20 April 2020?
\item Define \omterm{def:s1:calendar-spreads-and-curve-trades:storage}{storage-limit risk}.
\item Describe the rule's path through March and April 2020.
\item What did each filter do, and why is the best one suspect?
\end{enumerate}

\medskip\noindent\textbf{Part IV --- The verdict.}
\begin{enumerate}[resume]
\item State the \emph{named result}: the spread strategy's Sharpe ratio before 2020 and its loss in April 2020.
\item What would a defensible filter use?
\item What do the synthetic curves show about trading speed?
\item What did Simon and Campasano find about the VIX futures basis?
\item What is the volatility curve's equivalent of a storage limit?
\item Which strategy file depends on \omterm{def:s1:index-rebalancing:fund}{index funds}' rules?
\item How does this chapter's carry relate to chapter~20's?
\item In one sentence: what does a calendar spread hedge, and what does it not?
\end{enumerate}
\end{problem}

\section{Interview questions}

\begin{interviewq}[$\star$ \iqroles{researcher}]\label{iq:s1:calendar-spreads-and-curve-trades:1}
What determines the slope of a commodity futures curve?
\end{interviewq}

\begin{interviewq}[$\star\star$ \iqroles{developer}]\label{iq:s1:calendar-spreads-and-curve-trades:2}
How would you store futures data so that spread backtests follow contracts correctly through expiries?
\end{interviewq}

\begin{interviewq}[$\star\star$ \iqroles{trader}]\label{iq:s1:calendar-spreads-and-curve-trades:3}
You are long the WTI front spread and contango widens by 50\% a year in two days. What do you check, and what do you do?
\end{interviewq}

\begin{interviewq}[$\star\star$ \iqroles{risk}]\label{iq:s1:calendar-spreads-and-curve-trades:4}
How should a calendar spread book be margined and stress-tested?
\end{interviewq}

\begin{interviewq}[$\star\star$ \iqroles{researcher}]\label{iq:s1:calendar-spreads-and-curve-trades:5}
Why might short VIX futures in contango be profitable on average, and what is the risk?
\end{interviewq}

\begin{interviewq}[$\star\star\star$ \iqroles{researcher}]\label{iq:s1:calendar-spreads-and-curve-trades:6}
A log spread follows an AR(1) with coefficient $\phi$ and innovation standard deviation $\sigma$. A trader holds $-\,s_t$ units at the close of $t$.
Derive the expected daily P\&L and its Sharpe ratio, and show how a cost $k$ per unit traded changes it.
\end{interviewq}
